\section{Reproduction and evidence checklist}

A reviewer can begin with the application package scripts and the evaluation README. The production benchmark command is npm run benchmark:production, while npm run benchmark:comfort generates the separate restriction suite. The original journal-containing suite is retained for historical comparison. Strict ONNX mode must successfully load the configured model; a lexical run should remain separately labelled. Output files record engine version, backend identity, catalogue/\allowbreak{}scenario inputs, environment, seed, and source hashes. Repeating a command may overwrite that suite's standard output, so a custom output path should be used when preserving a saved result matters.

The relevant implementation paths are \path{recommender/app/engine.py} and schemas.py; \path{backend/src/routes/recommendations.ts} and \path{data/comfortPreferences.ts}; \path{src/services/recommendations/}; \path{src/services/journal/}; \path{src/context/JournalContext.tsx}; and \path{backend/src/routes/journal.ts}. Database definitions are in \path{backend/src/db.ts}. These are repository-relative paths and can change as the project evolves. Their inclusion provides a starting point for review rather than embedding source code or secrets into the dissertation.

The regression commands include npm run test:comfort, npm run test:journal, npm run recommender:test, and npm run recommender:typecheck. Mobile TypeScript, backend build, and web/\allowbreak{}native bundle exports check different integration boundaries. Journal API tests use synthetic fixtures and isolated database tooling; live native/\allowbreak{}backend tests require the intended services and installed app. A reviewer should record the date, source version, environment, commands, and results of any rerun. The saved passes cited in this report are attributed to their existing records rather than represented as new results.

The release checklist should verify the installed native version, API readiness, database migration status, recommendation model/\allowbreak{}version, storage policy, email delivery, provider client configuration, and journal recovery on an additional device. It should then exercise check-in, restriction enforcement, practice playback, encrypted reflection, attachment saving, history, logout, and deletion using synthetic accounts. Each item should have a recorded outcome and any limitation. No existing user's reflection should be used as an integration fixture.
